% paper/sections/08d_gfm.tex
%
% §8.4  Extension PDE による界面越し場延長
%
%   Aslam (2004) の Extension PDE により，界面 Γ を跨いで圧力場を
%   滑らかに延長し，CCD ステンシルの Gibbs 振動を除去する．
% ═══════════════════════════════════════════════════════════════════════════════
\subsection{Extension PDE による界面越し場延長}
\label{sec:gfm}
% ═══════════════════════════════════════════════════════════════════════════════

% NOTE: \label{sec:gfm} は §10 以降（検証実験）が参照するため変更不可．
% 節の内容は Extension PDE に完全に置き換えた（GFM は廃止）．

本節では，Aslam (2004)~\cite{Aslam2004} の \textbf{Extension PDE} を用いて
界面 $\Gamma$ を跨いで圧力場を滑らかに延長する手法を述べる（詳細は §~\ref{sec:extension_pde} 参照）．
CCD 法は同一ステンシル $(i-2,\ldots,i+2)$ を用いて $\Ord{h^6}$ の1・2階微分を
同時求解するが，界面を跨ぐ不連続場に適用すると Gibbs 振動を引き起こし精度が著しく低下する．
Extension PDE はこれを PPE 求解前に解消する前処理として機能する．

\subsubsection{Extension PDE の定式化}

スカラー場 $q$（圧力増分 $\delta p$ 等）を液相側（$\phi < 0$）から気相側（$\phi \ge 0$）へ滑らかに延長する
偏微分方程式は（付録~\ref{sec:extension_pde} も参照）：
\begin{equation}
  \frac{\partial q}{\partial \tau} + \mathrm{sgn}(\phi)\,\hat{\bm{n}}\cdot\bnabla q = 0
  \label{eq:extension_pde_main}
\end{equation}
ここで $\tau$ は仮想時間，$\hat{\bm{n}} = \bnabla\phi/|\bnabla\phi|$ は界面法線ベクトルであり，
$\mathrm{sgn}(\phi) = +1$（気相）$/-1$（液相）が延長方向を定める．
液相側（ソース相）の値を固定したまま，気相側（ターゲット相）を仮想時間前進させることで
$\nabla q \cdot \hat{\bm{n}} = 0$（法線方向勾配ゼロ）の定常解へ収束する．

\subsubsection{CCD との整合性}

式~\eqref{eq:extension_pde_main} を仮想時間前進により解くと，延長後の場 $q_{\mathrm{ext}}$ は
界面近傍で $C^\infty$ に近い滑らかな分布を持つ（詳細は §~\ref{sec:extension_pde} の図~\ref{fig:extension_schematic} 参照）．
この滑らかな場 $q_{\mathrm{ext}}$ に対して CCD 演算子 $D^{(1)}, D^{(2)}$ を適用することで：
\begin{itemize}
  \item 界面直近まで $\Ord{h^6}$ 精度の圧力勾配 $\bnabla(\delta p)$ が得られる
  \item 速度補正 $\bu^{n+1} = \bu^* - (\Delta t/\rho)\,\bnabla(\delta p)$ の離散化誤差が低減される
  \item Balanced-Force 条件（§~\ref{sec:balanced_force}）と整合した圧力勾配評価が実現される
\end{itemize}

\subsubsection{アルゴリズムへの組み込み}

時刻ループ内での適用箇所（§~\ref{sec:algorithm}）：
\begin{enumerate}
  \item \textbf{PPE 求解前（Step 5c）：} 前時刻圧力 $p^n$ に Extension PDE を適用して $p^n_{\mathrm{ext}}$ を構築し，
        IPC Predictor の $-\bnabla p^n$ 項に使用する
  \item \textbf{Corrector 前（Step 7 前）：} 圧力増分 $\delta p$ に Extension PDE を適用して $\delta p_{\mathrm{ext}}$ を構築し，
        速度補正 $\bu^{n+1} = \bu^* - (\Delta t/\rho)\,\bnabla(\delta p_{\mathrm{ext}})$ に使用する
\end{enumerate}

実装詳細（離散化，仮想時間刻み，イテレーション数）は §~\ref{sec:extension_pde} および
\texttt{twophase.levelset.field\_extender.FieldExtender} クラスを参照のこと．

% ─────────────────────────────────────────────────────────────────────────────
% Backward-compatibility labels for §10 cross-references (GFM equations removed)
% §10 以降が参照するため，ラベルのみ保持する（本文中の式は廃止）
\phantomsection\label{eq:gfm_rhs_correction}%
\phantomsection\label{eq:gfm_ccd_div}%
% ─────────────────────────────────────────────────────────────────────────────

% ═══════════════════════════════════════════════════════════════════════════════
